%% ICCE 2026 Sample Paper
%% Compile with: latexmk -xelatex -biber sample-paper.tex
%%
%% This sample demonstrates all ICCE formatting elements:
%% title, authors, abstract, keywords, sections, subsections,
%% figures (placeholder), tables, citations, acknowledgements, references.

\documentclass[11pt,a4paper]{article}

% ==== Font (XeLaTeX required) ====
\usepackage{fontspec}
\setmainfont{Arial}[
  BoldFont      = Arial Bold,
  ItalicFont    = Arial Italic,
  BoldItalicFont = Arial Bold Italic,
]

% ==== Page geometry ====
\usepackage[
  a4paper,
  top=1in,
  bottom=0.8in,
  left=1in,
  right=1in,
  nohead,
  nofoot,
]{geometry}

% ==== Line spacing ====
\usepackage{setspace}
\setstretch{1.0}

% ==== Paragraph indent ====
\setlength{\parindent}{1.25cm}
\setlength{\parskip}{0pt}

% ==== Section heading formatting ====
\usepackage{titlesec}
\titleformat{\section}
  {\normalfont\bfseries}{\thesection}{1em}{}
\titlespacing*{\section}
  {0pt}{2\baselineskip}{\baselineskip}

\titleformat{\subsection}
  {\normalfont\itshape}{\thesubsection}{1em}{}
\titlespacing*{\subsection}
  {0pt}{\baselineskip}{\baselineskip}

\titleformat{\subsubsection}
  {\normalfont\itshape}{\thesubsubsection}{1em}{}
\titlespacing*{\subsubsection}
  {0pt}{\baselineskip}{\baselineskip}

% ==== Title ====
\newcommand{\iccetitle}[1]{%
  \begin{center}
    {\fontsize{22pt}{26pt}\selectfont\bfseries #1}%
  \end{center}%
  \vspace{\baselineskip}%
}

% ==== Author ====
\newcommand{\icceauthor}[2]{%
  \begin{center}
    #1 \\ {\itshape #2}%
  \end{center}%
}

% ==== Abstract ====
\newenvironment{icceabstract}{%
  \begin{list}{}{%
    \setlength{\leftmargin}{0.5in}%
    \setlength{\rightmargin}{0.5in}%
  }%
  \item[]\small\noindent\textbf{Abstract:}\space
}{%
  \end{list}%
}

% ==== Keywords ====
\newenvironment{iccekeywords}{%
  \vspace{\baselineskip}%
  \begin{list}{}{%
    \setlength{\leftmargin}{0.5in}%
    \setlength{\rightmargin}{0.5in}%
  }%
  \item[]\small\noindent\textbf{Keywords:}\space
}{%
  \end{list}%
  \vspace{2\baselineskip}%
}

% ==== Figures and tables ====
\usepackage{graphicx}
\usepackage{caption}
\usepackage{booktabs}
\captionsetup[figure]{justification=centering, position=below}
\captionsetup[table]{justification=centering, position=above}

% ==== Bibliography (APA 7th) ====
\usepackage[style=apa, backend=biber, natbib=true]{biblatex}
\addbibresource{sample-paper.bib}

% ==== Lists ====
\usepackage{enumitem}

% ==== Bookmarks and accessibility ====
\usepackage[
  pdfauthor={Shota MIURA, Masanori TAKAGI},
  pdftitle={Generating Programming Quizzes with Retrieval-Augmented Generation: An Experimental Study with High School Students},
  bookmarks=true,
  bookmarksopen=true,
  colorlinks=true,
  linkcolor=black,
  citecolor=black,
  urlcolor=blue,
]{hyperref}
\usepackage{bookmark}

% ==== No page numbers ====
\pagestyle{empty}

% ================================================================
\begin{document}

\vspace{\baselineskip}

\iccetitle{Generating Programming Quizzes with Retrieval-Augmented Generation: An Experimental Study with High School Students}

\icceauthor{Shota MIURA, Masanori TAKAGI}%
  {The University of Electro-Communications, Chofu, Tokyo, Japan}

\vspace{\baselineskip}

\begin{icceabstract}
Creating programming quizzes that match learners' proficiency levels is a time-consuming task for instructors. This study proposes a method for generating programming quizzes using Retrieval-Augmented Generation (RAG) with large language models (LLMs). The proposed system retrieves relevant educational materials from a curated knowledge base and uses them as context for quiz generation, aiming to produce questions that are both technically accurate and pedagogically appropriate. We conducted an experiment with 32 high school students at a public senior high school in the western suburbs of Tokyo, Japan. The students were in the 11th grade (the second year of a three-year senior high school curriculum in Japan, ages 16--17). Results showed that quizzes generated by the RAG-based system received significantly higher ratings for appropriateness and clarity compared to quizzes generated without retrieval augmentation. The findings suggest that incorporating domain-specific retrieval into LLM-based quiz generation can improve the educational quality of automatically generated questions.
\end{icceabstract}

\begin{iccekeywords}
retrieval-augmented generation, programming education, quiz generation, large language models, high school students
\end{iccekeywords}

% ================================================================

\section{Introduction}

Programming education in secondary schools has expanded rapidly in recent years, particularly in Japan where information science became a required subject in the 2022 national curriculum revision. A key challenge for instructors is creating assessment materials---such as quizzes and exercises---that are aligned with both the curriculum and the varied proficiency levels of learners \parencite{smith2024}. Traditionally, this task has required significant manual effort from teachers who must design questions, verify their correctness, and ensure appropriate difficulty levels.

Recent advances in large language models (LLMs) offer promising opportunities for automating quiz generation. However, LLMs trained on general corpora may produce questions that are misaligned with specific curricula or that contain technical inaccuracies. Retrieval-Augmented Generation (RAG) addresses this limitation by augmenting the generation process with relevant retrieved documents, grounding the output in domain-specific knowledge \parencite{wang2024}.

In this study, we propose a RAG-based system for generating programming quizzes and evaluate its effectiveness through an experimental study with high school students. The research questions are as follows:

\begin{enumerate}
  \item Can RAG-based quiz generation produce questions that are more appropriate and clearer than those generated without retrieval augmentation?
  \item How do learners perceive the difficulty and educational value of RAG-generated quizzes?
\end{enumerate}

\section{Related Work}

\subsection{Automated Quiz Generation}

Automated quiz generation has been studied extensively in educational technology. Early approaches relied on template-based methods that filled predefined question structures with content extracted from textbooks. More recently, neural approaches using transformer-based models have shown improved fluency and diversity in generated questions \parencite{chen2023}. However, these methods often lack domain-specific grounding, leading to factual errors or misalignment with learning objectives.

\subsection{Retrieval-Augmented Generation in Education}

RAG combines the generative capabilities of LLMs with the precision of information retrieval systems. In educational contexts, RAG has been applied to personalized learning content generation \parencite{wang2024}, intelligent tutoring systems, and adaptive assessment. The key advantage is that retrieved documents provide factual grounding, reducing hallucinations and improving the relevance of generated content.

\section{Method}

\subsection{System Design}

The proposed system consists of three main components: (1) a knowledge base of programming educational materials, (2) a retrieval module that selects relevant documents based on topic and difficulty, and (3) a generation module that produces quiz questions using the retrieved context.

\subsection{Quiz Generation Process}

Given a topic and target difficulty level, the system first retrieves the top-$k$ most relevant documents from the knowledge base using semantic similarity search. These documents are then provided as context to the LLM, which generates quiz items following a structured prompt template. The generated quizzes include multiple-choice questions, code completion tasks, and short-answer questions.

% FIGURE HERE: System architecture diagram showing the three components and data flow

\section{Experimental Application}

\subsection{Participants}

The participants were 32 high school students at a public senior high school in the western suburbs of Tokyo, Japan. The students were in the 11th grade (the second year of a three-year senior high school curriculum in Japan, ages 16--17). All participants had completed an introductory Python programming course covering variables, conditionals, and loops.

\subsection{Procedure}

The experiment followed a within-subjects design. Each participant completed two sets of programming quizzes: one generated by the RAG-based system and one generated by the same LLM without retrieval augmentation. The order of quiz sets was counterbalanced across participants. After completing each quiz, participants rated the following on a 5-point Likert scale:

\begin{itemize}
  \item Appropriateness of question content
  \item Clarity of question wording
  \item Perceived difficulty
  \item Educational value
\end{itemize}

\subsection{Results}

Table~\ref{tab:ratings} summarizes the average ratings for each condition.

\begin{table}[htbp]
  \centering
  \caption{Average Ratings on a 5-Point Likert Scale}
  \label{tab:ratings}
  \begin{tabular}{lcc}
    \toprule
    Dimension & RAG-based & Without RAG \\
    \midrule
    Appropriateness & 4.21 & 3.45 \\
    Clarity         & 4.08 & 3.52 \\
    Difficulty      & 3.56 & 3.61 \\
    Educational value & 4.15 & 3.38 \\
    \bottomrule
  \end{tabular}
\end{table}

A paired-samples $t$-test revealed that the RAG-based condition received significantly higher ratings for appropriateness ($t(31) = 4.12$, $p < .001$), clarity ($t(31) = 3.47$, $p < .01$), and educational value ($t(31) = 3.89$, $p < .001$). No significant difference was found for perceived difficulty ($t(31) = 0.34$, $p = .74$).

% FIGURE HERE: Bar chart comparing ratings between RAG and non-RAG conditions

\section{Analysis and Discussion}

The results indicate that incorporating retrieval augmentation into LLM-based quiz generation significantly improves the appropriateness, clarity, and educational value of generated quizzes. This finding aligns with prior work showing that RAG reduces hallucinations and improves factual accuracy in generated content \parencite{wang2024}.

The absence of a significant difference in perceived difficulty suggests that retrieval augmentation does not systematically alter the difficulty level of generated questions. This is desirable, as difficulty should be controlled by the instructor through parameter settings rather than being an artifact of the generation method.

One limitation of this study is the relatively small sample size ($n = 32$). Future work should replicate these findings with a larger and more diverse participant pool. Additionally, the knowledge base used in this study was curated for Python programming; extending the system to other domains would require building domain-specific knowledge bases.

\section{Conclusion}

This study proposed a RAG-based system for generating programming quizzes and evaluated it through an experiment with high school students. The results demonstrated that RAG-based quiz generation produces questions that are significantly more appropriate, clearer, and educationally valuable than those generated without retrieval augmentation. These findings contribute to the growing body of research on applying LLMs to educational assessment and suggest that retrieval augmentation is a promising approach for improving the quality of automatically generated educational content.

% ==== Acknowledgements (unnumbered) ====
\section*{Acknowledgements}
\addcontentsline{toc}{section}{Acknowledgements}

This work was supported by the Japan Society for the Promotion of Science (JSPS) KAKENHI Grant Number XXXXXXX. We thank the teachers and students at the participating high school for their cooperation.

% ==== References ====
\printbibliography[title={References}]

\end{document}
